\section{迹态射}\etiquette[2]{XVIII.sec2}\pageoriginale

本 \S 致力于证明定理 \hyperref[XVIII.thm2.9]{2.9} 及 \hyperref[XVIII.prop2.12]{2.12}. \hyperref[XVIII.thm2.9]{2.9} 的证明仅使用第 \hyperref[XVIII.sec1]{1} \S 中的第 1 小节.

\begin{enonce*}{引理 2.1}\etiquette[2.1]{XVIII.lem2.1}
设 $f:X\rightarrow S$ 为相对维数 $\leq d$ 的可紧化态射 \emph{(\hyperrefext[XVII][XVII.def3.2.1]{3.2.1})}, $j:U\hookrightarrow X$ 为 $X$ 的有限型开子概形, 其补在每个纤维中的维数 $< d$ (例如, $U$ 在 $X$ 的每个纤维中稠密). 若 $F$ 为 $X$ 上的挠层, 则典范箭头
$$
R^{2d}(fj)_{!}j^{\ast}F\xrightarrow{\sim} R^{2d}f_{!}F
$$
为同构.
\end{enonce*}

设 $i$ 为 $Y=X-U$ 到 $X$ 的嵌入. 正合列 (XVII \hyperref[XVII.eq5.1.16.2]{5.1.16.2}) 给出正合列
$$
R^{2d-1}(fi)_{!}i^{\ast}F\longrightarrow
R^{2d}(fj)_{!}j^{\ast}F\longrightarrow R^{2d}f_{!}F\longrightarrow
R^{2d}(fi)_{!}i^{\ast}F
$$
由 (\hyperrefext[XVII][XVII.cor5.2.8.1]{5.2.8.1}), 其两端的项为零.

\begin{enonce*}{引理 2.2}\etiquette[2.2]{XVIII.lem2.2}
设 $f:X\rightarrow S$ 为相对维数 $\leq d$ 的可紧化态射, $u_{i}:X_{i}\rightarrow X$ 为有限型分离平展态射族; 置 $u_{ij}:X_{ij}=\fprod{X_{i}}{X_{j}}{X}\rightarrow X$. 对 $X$ 上的每个挠层 $F$, 箭头 \emph{(XVII \hyperref[XVII.eq6.2.7.2]{6.2.7.2})} 所成的序列
$$
\bigoplus_{i,j} R^{2d}(fu_{ij})_{!}u^{\ast}_{ij}F\rightrightarrows
\bigoplus_{i} R^{2d}(fu_{i})_{!}u^{\ast}_{i}F\rightarrow
R^{2d}f_{!}F\rightarrow 0
$$
正合.
\end{enonce*}

这立即由函子 $R^{2d}f_{!}$ 的右正合性及 \hyperrefext[XVII][XVII.prop6.2.9]{6.2.9} 得到.

\begin{enonce*}{引理 2.3}\etiquette[2.3]{XVIII.lem2.3}
设 $f:X\rightarrow S$ 为有限表示的平坦态射. $X$ 中 $f$ 为 Cohen--Macaulay 的开子集 \emph{(\EGA IV 6.8.1 及 12.2.1 (vii))} 相对有限表示, 且在每个纤维中稠密.
\end{enonce*}

此\pageoriginale 问题在 $S$ 上是局部的, 归约后可假设其为 Noether 概形, 从而得到第一个断言. 第二个断言逐纤维验证, 因而可取 $S$ 为域的谱. $X$ 在 $X$ 的极大点处的局部环是 Artin 环, 因而为 Cohen--Macaulay 环; 由于“Cohen--Macaulay”是“开”性质, 断言得证.

\begin{enonce*}{引理 2.4}\etiquette[2.4]{XVIII.lem2.4}
设 $f:X\rightarrow S$ 为有限表示、平坦且 Cohen--Macaulay 的态射, $S$ 拟分离. $X$ 的每个点都有一个 (Zariski) 开邻域 $U$, 使存在从 $U$ 到标准向量丛 $\EE^{d}_{S}$ 的拟有限、平坦且有限表示的 $S$ 态射.
\end{enonce*}

这是 \EGA IV 15.4.3 的特殊情形.

\begin{enonce*}{引理 2.5}\etiquette[2.5]{XVIII.lem2.5}
设 $\underline{u}=(u_{1},\ldots,u_{d})$ 与 $\underline{v}=(v_{1},\ldots,v_{d})$ 为维数为 $d$、极大理想为 $\mathfrak{m}$ 的 Cohen--Macaulay 局部环 $A$ 的两组参数系. 存在 $A$ 的参数系序列 $\underline{w}_{0},\ldots,\underline{w}_{n}$, 使 $\underline{w}_{0}=\underline{u}$、$\underline{w}_{n}=\underline{v}$, 且每个 $\underline{w}_{i+1}$ 只改变 $\underline{w}_{i}$ 中的一个参数而得.
\end{enonce*}

对 $d$ 作归纳, 当 $d<2$ 时断言为空泛的; 因而假设 $d\geq 2$. 元素 $x\in \mathfrak{m}$ 为 $A/u_{1}$ 正则 (\resp $A/v_{1}$ 正则) 的充要条件是 $x$ 不属于 $\Ass(A/u_{1})$ (\resp $\Ass(A/v_{1})$) 的任何成员; 由于素理想 $\neq \mathfrak{m}$ 的任何有限并都不等于 $\mathfrak{m}$, 存在同时为 $A/u_{1}$ 正则与 $A/v_{1}$ 正则的 $x\in \mathfrak{m}$, 因而也存在参数系 $\underline{a}=(u_{1},x,\ldots)$ 与 $\underline{b}=(v_{1},x,\ldots)$. 在 $A/u_{1}$ 中对 $\underline{u}-\{u_{1}\}$ 与 $\underline{a}-\{u_{1}\}$、在 $A/x$ 中对 $\underline{a}-\{x\}$ 与 $\underline{b}-\{x\}$、在 $A/v_{1}$ 中对 $\underline{b}-\{v_{1}\}$ 与 $\underline{v}-\{v_{1}\}$ 应用归纳假设, 即得结论.

\begin{enonce*}{引理 2.6}\etiquette[2.6]{XVIII.lem2.6}
设 $X$ 为代数闭域 $k$ 上的有限型 Cohen--Macaulay 概形, $u$、$v:X\rightarrow \EE^{d}_{k}$ 为从 $X$ 到标准向量丛的两个拟有限平坦态射. $X$ 的每个闭点 $x$ 都有一个 (Zariski) 开邻域 $U$, 使存在态射序列\pageoriginale $w_{0},\ldots,w_{n}:U\rightarrow \EE^{d}_{k}$, 满足 $w_{0}=u|U$、$w_{n}=v|U$, 且 $w_{i+1}$ 与 $w_{i}$ 仅有一个坐标不同.
\end{enonce*}

先假设 $u(x)=v(x)=0$. 此时, 在 $x$ 处从 $X$ 到 $\AAA^{d}_{k}$ 的拟有限平坦态射的芽,

若将 $x$ 映到 $0$, 就等同于局部环 $\Oo_{X,x}$ 的参数系, 而 \hyperref[XVIII.lem2.6]{2.6} 由 \hyperref[XVIII.lem2.5]{2.5} 得到. 在一般情形下, 先用态射链将 $u$ 接到 $u-u(x)$ (\resp 将 $v$ 接到 $v-v(x)$):
$$
u-(u_{1}(x),\ldots,u_{i}(x),0,\ldots,0)(\text{\resp}
v-(v_{1}(x),\ldots,v_{i}(x),0,\ldots,0)){\quoi}.
$$

\setcounter{subsection}{6}
\subsection{}\etiquette[2.7]{XVIII.subsec2.7}
设有一对可复合的 $S$ 可紧化态射 $X\xrightarrow{g} Y\xrightarrow{f}S$. 假设 $f$ 的相对维数 $\leq d$, $g$ 的相对维数 $\leq
e$. 则对 $X$ 上的每个挠层 $F$, 复合谱序列
\begin{equation*}
R^{p}f_{!}R^{q}g_{!}F\Longrightarrow R^{p+q}(fg)_{!}F\tag{2.7.1}
\end{equation*}\etiquette[2.7.1]{XVIII.eq2.7.1}%
给出同构 (利用最高次循环项的论证)
\begin{equation*}
R^{2d}f_{!}R^{2e}g_{!}F\simeq R^{2(d+e)}(fg)_{!}F{\quoi}.\tag{2.7.2}
\end{equation*}\etiquette[2.7.2]{XVIII.eq2.7.2}

\subsection{}\etiquette[2.8]{XVIII.subsec2.8}
设
$$
a^{d}:\EE^{d}_{S}\longrightarrow S
$$
为凝聚 (即拟紧、拟分离) 概形 $S$ 上的标准向量丛. 若 $F$ 为 $S$ 上与 $S$ 的剩余特征互素的挠层, 则如下对 $d$ 作归纳, 定义迹态射
\begin{equation*}
\Tr_{a^d}:R^{2d}a^{d}_{!}F(d)\xrightarrow{\sim} F{\quoi}.\tag{2.8.1}
\end{equation*}\etiquette[2.8.1]{XVIII.eq2.8.1}

当 $d=0$ 时, $\Tr_{a^0}$ 为恒等态射. 当 $d=1$ 时, $\Tr_{a^1}$ 为同构 (\hyperref[XVIII.prop1.1.6]{1.1.6}) (\cf \hyperref[XVIII.lem1.1.9]{1.1.9}). 最后, 有典范同构
$$
\EE^{d+1}_{S}=\fprod{\EE^{1}_{S}}{\EE^{d}_{S}}{S}=\EE^{1}_{\EE_{S}^d}
$$
将\pageoriginale $\Tr_{a^{d+1}}$ 定义为同构 $\Tr_{a^{d}_{S}}$ 与
$Ra^{d}_{S!}(\Tr_{a^1_{\EE^d_S}})$
的复合 (\cf \hyperref[XVIII.thm2.9]{2.9} 的图表 Var 3).

\eqhyperref[XVIII.eq2.8.1]{2.8.1} 的源 (及靶) 的形成与任意基变换相容. 因而每个在 $\EE^{d}_{S}$ 上作用的 $S$ 上的代数群, 都作用于这些层; 连通代数群的作用必为平凡作用. 因此, 包含于仿射群中的坐标置换群在 \eqhyperref[XVIII.eq2.8.1]{2.8.1} 的源 (及靶) 上的作用是平凡的.

\setcounter{subsubsection}{1}
\subsubsection{}\etiquette[2.8.2]{XVIII.subsubsec2.8.2}
同构 \eqhyperref[XVIII.eq2.8.1]{2.8.1} 在坐标置换下不变.

也可用杯积解释 \eqhyperref[XVIII.eq2.8.1]{2.8.1}, 从而验证这一点 (XVII \hyperref[XVII.eq5.4.2.2]{5.4.2.2} 及 \hyperrefext[XVII][XVII.lem5.4.3.5]{5.4.3.5}).

\begin{enonce*}{定理 2.9}\etiquette[2.9]{XVIII.thm2.9}
考虑三元组 $(f,d,F)$, 由可紧化态射 \emph{(\hyperrefext[XVII][XVII.subsec6.3]{6.3})} $f:X\rightarrow Y$、整数 $d$ 及 $Y$ 上与 $Y$ 的剩余特征互素的挠层 $F$ 组成, 其中态射 $f$ 满足条件 $(\ast)_{d}$: 存在 $X$ 的开子集 $U$, 使 $f$ 在 $U$ 上的限制为有限表示的平坦态射, 其纤维维数为 $\leq d$, 并使 $X-U$ 的纤维维数 $< d$.

有且仅有一种方式, 对每个上述三元组 $(f,d,F)$ 关联迹态射
$$
\Tr_{f}:R^{2d}f_{!}f^{\ast}F(d)\longrightarrow F
$$
使下列条件满足:
\begin{itemize}
\item[(Var 1)] (函子性). 迹态射对 $F$ 具有函子性.

\item[(Var 2)] (与基变换相容). 对每个笛卡尔方块
$$
\xymatrix{
X' \ar[r]_{g'}\ar[d]^{f'}& X\ar[d]^{f}\\
Y' \ar[r]_{g} & Y
}
$$
其中\pageoriginale $Y'$ 拟紧且拟分离, 图表
$$
\xymatrix@R=1.25cm{
R^{2d}f'_!f'^{\ast}g^{\ast}F(d) \ar@{=}[r]\ar[d]^{Tr_{f'}}& R^{2d}f'_!
g'^{\ast} f^{\ast} F(d) & g^{\ast}
R^{2d}f_!f^{\ast} F(d) \ar[l]_{\sim} \ar[d]^{g^{\ast}(Tr_f)}\\
g^{\ast}F \ar@{=}[rr]&            & g^{\ast}F
}
$$
交换, 其中水平同构为基变换箭头 \emph{\eqhyperrefext[XVII][XVII.eq5.2.2.1]{5.2.2.1}}.

\item[(Var 3)] (与复合相容). 对每一对可复合的 $Z$ 可紧化态射
$$
X\xrightarrow{g}Y\xrightarrow{f}Z{\quoi},
$$
以及使 $f$ 满足 $(\ast)_{d}$、$g$ 满足 $(\ast)_{e}$ 的整数 $d$ 与 $e$, 和 $Z$ 上的挠层 $F$, 复合态射 $fg$ 满足 $(\ast)_{d+e}$, 且图表
$$
\xymatrix@C=2.25cm{
R^{2d}f_! R^{2e} g_! g^{\ast} f^{\ast}F(d)(e) \ar@{<->}[d]_{\text{\rotatebox{90}{$\backsim$}}} \ar[r]^(.6){R^{2d}f_!(\Tr_g)} &
R^{2d}f_! f^{\ast}F(d) \ar[d]^{\Tr_f}\\
R^{2(d+e)} (fg)_! (fg)^{\ast} F(d+e) \ar[r]^(.65){\Tr_{fg}} & F
}
$$
交换, 其中竖直同构为箭头 \eqhyperref[XVIII.eq2.7.2]{2.7.2}.

\item[(Var 4)] (规范化条件).
\begin{itemize}
\item[(I)] 若 $f$ 为秩为 $n$ 的有限局部自由态射, 且 $d=0$, 则复合态射
$$
F\longrightarrow
f_{\ast}f^{\ast}F=f_{!}f^{\ast}F\xrightarrow{\tr_{f}}F
$$
为乘以 $n$.

\item[(II)] 若\pageoriginale $f$ 为标准仿射直线 $f:\EE^{1}_{Y}\rightarrow Y$, 且 $d=1$, 则迹态射
$$
\Tr_{f}: R^{2}f_{!}f^{\ast}F(1)\longrightarrow F
$$
为同构 \eqhyperref[XVIII.eq2.8.1]{2.8.1}.
\end{itemize}
\end{itemize}
\end{enonce*}

\begin{enumerate}
\renewcommand{\theenumi}{\alph{enumi}}
\renewcommand{\labelenumi}{\theenumi)}
\item 对于 $d=0$, 条件 $(\ast)_{d}$ 表示 $f$ 拟有限、平坦且有限表示. 由 (\hyperrefext[XVII][XVII.thm6.2.3]{6.2.3}), 对 $d=0$, 迹态射必与 \loccit{} 中所构造的态射一致.

当 $f$ 为标准仿射空间 $\AAA^{d}_{S}$ 到 $S$ 的投影时, 由 (Var~4)~(II) 及 (Var~3) 可知, $\Tr_{f}$ 必须为同构 \eqhyperref[XVIII.eq2.8.1]{2.8.1}.

\item 设 $f:X\rightarrow Y$ 为满足条件 $(\ast)_{d}$ 的可紧化态射: 存在从 $X$ 到标准向量丛 $a^{d}:\EE^{d}_{Y}\rightarrow Y$ 的拟有限、平坦且有限表示的 $Y$ 态射 $u$.

设 $t(f,u)$ 为迹态射 \eqhyperref[XVIII.eq2.8.1]{2.8.1} 与 (\SGA \hyperrefext[XVII][XVII.thm6.2.3]{6.2.3}) 的“复合”态射 (\cf (Var 3))
$$
t(f,u)=\Tr_{a^d}\circ R^{2d}a^{d}_{!}(\Tr_{u}){\quoi}.
$$

此态射 $t(f,u)$ 满足 (Var 1) 及 (Var 2). 通过基变换, 为验证 $t(f,u)$ 只依赖于 $f$ 及 $d$, 而不依赖于 $u$, 只须在 $Y$ 为代数闭域 $k$ 的谱时验证之. 由 \hyperref[XVIII.lem2.2]{2.2}, 此问题在 $X$ 上是局部的, 所以由 \hyperref[XVIII.lem2.6]{2.6}, 只须证明
$t(f,u)=t(f,v)$
其中 $u$ 与 $v$ 只有一个坐标不同; 由 \hyperref[XVIII.subsubsec2.8.2]{2.8.2}, 可以假设此坐标为第一个. 置 $E=\EE^{d-1}_{Y}$, 于是有交换图表
$$
\xymatrix{
                              &    &\EE^1_{E} \ar[rdd]_(.3){a} &  &  \\
X \ar[rrrr]_{v}\ar[rru]^{u} \ar[rdd]_{f} \ar[rrrd]_{x}&    &                 &
                              &\EE^{1}_E \ar[ld]_{a}\\
                              &    &                 &E \ar[lld]^w& \\
                              & Y  &                 &  &
}
$$
根据\pageoriginale 构造 (\hyperref[XVIII.subsec2.8]{2.8}),
$$
\Tr_{wa}=\Tr_{w}\circ R^{2(d-1)}w_!(\Tr_{a}){\quoi}\;;
$$
推得 $t(f,u)$ 为 $\Tr_{w}$ 与 $t(x,u)$ 的“复合”:
\begin{align*}
t(f,u) &= \Tr_{wa}\circ R^{2d}(wa)_{!}(\Tr_{u})\\
       &= \Tr_{w}\circ R^{2(d-1)}w_{!}(\Tr_{a})\circ
       R^{2d}(wa)_{!}(\Tr_{u})\\
       &= \Tr_{w}\circ R^{2(d-1)}w_{!}(\Tr_{a}\circ
       R^{2}a_{!}(\Tr_{u})){\quoi}.
\end{align*}
此外 (\hyperref[XVIII.lem1.1.7]{1.1.7}), $t(x,u)$ 恰为 \hyperref[XVIII.prop1.1.6]{1.1.6} 的态射 $\Tr_{x}$; 因而 $t(x,u)=t(x,v)=\Tr_{x}$ 且 $t(f,u)=t(f,v)$.

以上面构造的态射 $t(f,u)$ 记为 $\Tr_{f}$. 由 \hyperrefext[XVII][XVII.thm6.2.3]{6.2.3}, 此迹态射满足 (Var 1)、(Var 2)、$e=0$ 时的 (Var 3) 及 (Var 4). 只要它有定义, 就必与所求迹态射一致, 如果后者存在.

\item 设 $f:X\rightarrow Y$ 为可紧化的 Cohen--Macaulay 态射, 平坦且有限表示, 纯相对维数为 $d$. 设 $(U_{i})_{i\in I}$ 为 $X$ 的拟紧开子集覆盖, 使存在从 $U_{i}$ 到 $\EE^{d}_{Y}$ 的拟有限平坦 $Y$ 态射 (\hyperref[XVIII.lem2.4]{2.4}). 由 \hyperref[XVIII.lem2.2]{2.2} 及 b), 存在唯一态射 $\Tr_{f}$ 使下图交换, 这里采用 \hyperref[XVIII.lem2.2]{2.2} 的记号:
$$
\xymatrix{
\bigoplus_{i,j}
R^{2d}(fu_{ij})_!(fu_{ij})^{\ast} F(d) \ar@<2pt>[r]
\ar@<-2pt>[r]&\bigoplus_{i}
R^{2d}(fu_i)_!(fu_i)^{\ast}F(d)
\ar[rdd]_{\bigoplus_{i}\Tr_{fu_i}}\ar[r]
&
R^{2d}f_!f^{\ast}F(d) \ar[dd]^{\Tr_f}\\
& & \\
& & F
}
$$

此迹态射不依赖于所选取的开覆盖 (用并覆盖比较两个覆盖); 它满足 (Var 1)、(Var 2)、$e=0$ 时的 (Var 3) 及 (Var 4), 这立即由 b) 得到.

\item 设 $f:X\rightarrow Y$ 为满足 $(\ast)_{d}$ 的可紧化态射, $j:U\hookrightarrow X$\pageoriginale 为 $(\ast)_{d}$ 所保证的开子集. 根据 \hyperref[XVIII.lem2.3]{2.3}, 将 $U$ 换成较小开子集后, 可以假设 $fj$ 为 Cohen--Macaulay 且纯维数为 $d$.

\noindent
由 \hyperref[XVIII.lem2.1]{2.1}, 典范态射
$$
R^{2d}(fj)_{!}(fj)^{\ast}F(d)\xrightarrow{\sim}
R^{2d}f_{!}f^{\ast}F(d)
$$
为同构, 将 $\Tr_{f}$ 定义为复合
$$
\Tr_{f}:R^{2d}f_{!}f^{\ast}F(d)\xleftarrow{\sim}
R^{2d}(fj)_{!}(fj)^{\ast}F(d)\xrightarrow{\Tr_{fj}}F{\quoi}.
$$
由 c) 中建立的 $e=0$ 时的 (Var 3), 立即可知 $\Tr_{f}$ 与 $U$ 的选择无关. 此迹态射满足 (Var 1)、(Var 2)、$e=0$ 时的 (Var 3) 及 (Var 4), 且是唯一能满足所加条件的态射. 还须证明它满足 (Var 3), 为此, 上述论证将我们归结为只考虑分别满足 $(\ast)_{d}$ 与 $(\ast)_{e}$ 的对 $(f,d)$、$(g,e)$.

希望细致整理以下论证的读者, 可先验证下述引理:

\begin{enonce*}{引理 2.9.1}\etiquette[2.9.1]{XVIII.lem2.9.1}
设 $(f,g,h)$ 为三个可复合的 $T$ 可紧化态射, 满足 $(\ast)_{k}$、$(\ast)_{\ell}$ 及 $(\ast)_{m}$.
$$
X\xrightarrow{h}Y\xrightarrow{g}Z\xrightarrow{f}T{\quoi}.
$$
若对 $(f,g)$ 与 $(g,h)$ 满足 {\rm (Var 3)}, 则 $(f,gh)$ 满足 {\rm (Var 3)} 当且仅当 $(fg,h)$ 满足 {\rm (Var 3)}.
\end{enonce*}

结论中的条件也表示, $\Tr_{fgh}$ 为 $\Tr_{f}$、$\Tr_{g}$ 及 $\Tr_{h}$ 的“复合”.

\item 设给定交换图表:
$$
\xymatrix@R=1.25cm@C=1.5cm{
X \ar[r]^{v} \ar[rd]_{g}& \EE^{e}_Y \ar[r]^{u'}\ar[d]^{c} & \EE_Z^{d+e}\ar[d]^{b} \\
            &     Y   \ar[r]^{u} \ar[rd]_{f}                       & \EE_Z^{d}\ar[d]^{a}  \\
            &                                  &     Z
}
$$
相对于\pageoriginale $fg$ 的迹态射为相对于 $ab$ 与 $u'v$ 的迹态射的“复合”; 根据 $ab$ 时的构造以及 $u'v$ 时的 \hyperrefext[XVII][XVII.thm6.2.3]{6.2.3}, 它也是相对于 $a$、$b$、$u'$ 及 $v$ 的迹态射的“复合”. 另一方面, 相对于 $f$ 与 $g$ 的迹态射的“复合”, 是相对于 $a$、$u$、$c$ 及 $v$ 的迹态射的“复合”. 因而必须验证, 对每个笛卡尔图表
\begin{equation*}
\vcenter{
\xymatrix@C=1.25cm@R=1cm{
\EE^{e}_S  \ar[d]^{b'} \ar[r]^{u'}   & \EE^{e}_T \ar[d]^{b}\\
    S      \ar[r]^{u}     &      T
}}\tag{2.9.2}
\end{equation*}\etiquette[2.9.2]{XVIII.eq2.9.2}%
其中 $u$ 拟有限、平坦且有限表示, 迹态射 \eqhyperref[XVIII.eq2.8.1]{2.8.1} 及 \hyperrefext[XVII][XVII.subsec6.2]{6.2} 的图表
\begin{equation*}
\vcenter{
\xymatrix@C=2.25cm{
R^{2e}b_!u'_! u'^{\ast}b {}^{\ast} F(e) \ar[r]^(.55){R^{2e}b_!
  (\Tr_{u'})}\ar@{=}[d]& R^{2e}b_! b^{\ast} F(e)\ar[r]^{\Tr_b}
&  F \ar@{=}[d]\\
u_! R^{2e}b'_! b'^{\ast} u^{\ast} F(e)\ar[r]^(.55){u_!(\Tr_{b'})} & u_!
u^{\ast} F \ar[r]^{\Tr_u}& F
}}\tag{2.9.3}
\end{equation*}\etiquette[2.9.3]{XVIII.eq2.9.3}%
交换.
\end{enumerate}


通过基变换, 只须在 $T$ 为代数闭域 $k$ 的谱时验证之. 若 $S$ 为若干概形 $S_{i}$ 的和, 只须对以 $u_{i}:S_{i}\rightarrow T$ 为基的图表 \eqhyperref[XVIII.eq2.9.2]{2.9.2} 验证 \eqhyperref[XVIII.eq2.9.3]{2.9.3} 的交换性. 因而可以假设 $S$ 为有限次数 $n=[A:k]$ 的局部 Artin $k$ 代数 $A$ 的谱. 此时若通过函子 $u^{\ast}$ (\resp ${u'}^{\ast}$) (VIII \hyperref[VIII.thm1.1]{1.1}) 将 $T$ (\resp $\EE^{e}_{T}$) 上的层等同于 $S$ (\resp $\EE^{e}_{S}$) 上的层, 则态射 $\Tr_{b}$ 等同于态射 $\Tr_{b'}$ (\eqhyperref[XVIII.eq2.8.1]{2.8.1} 与基变换的相容性), 而 $\Tr_{u}$ 与 $\Tr_{u'}$ 等同于乘以 $n$, 断言得证\nde{也可将此相容性归结为情形 $e=1$, 它由结合 \hyperref[XVIII.lem1.1.7]{1.1.7} 及 \hyperref[XVIII.lem1.1.8]{1.1.8} 得到. 此外, 这里所给的论证与 \hyperref[XVIII.lem1.1.8]{1.1.8} 中默然略过的论证类似.}. 至此完成了 \hyperref[XVIII.thm2.9]{2.9} 的证明.

在证明过程中已经\pageoriginale 看到:

\begin{enonce*}{命题 2.10}\etiquette[2.10]{XVIII.prop2.10}
迹态射 \emph{\hyperrefext[XVII][XVII.thm6.2.3]{6.2.3}、XVIII~\hyperref[XVIII.prop1.1.6]{1.1.6}} 及 \eqhyperref[XVIII.eq2.8.1]{2.8.1} 均为迹态射 \hyperref[XVIII.thm2.9]{2.9} 的特殊情形.
\end{enonce*}

\begin{enonce*}[remark]{注 2.10.1}\etiquette[2.10.1]{XVIII.rem2.10.1}
不难证明, \hyperref[XVIII.thm2.9]{2.9} 的迹态射为同构, 当且仅当 $f$ 的几何纤维恰有一个维数为 $d$ 的不可约分支, 且其“重数”与 $n$ 互素. 这里不给出这一事实的直接证明, 它将由第 \hyperref[XVIII.sec3]{3} \S 的“整体对偶定理”立即得到.
\end{enonce*}

\setcounter{subsection}{10}
\subsection{}\etiquette[2.11]{XVIII.subsec2.11}
设 $f:X\rightarrow S$ 与 $g:Y\rightarrow S$ 为两个 $S$ 可紧化态射, 分别满足 $(\ast)_{d}$ 及 $(\ast)_{e}$.
$$
\xymatrix{
     & \fprod{X}{Y}{S} \ar[dd]^{f\times g}\ar[ld]_{\pr_1}\ar[rd]^{\pr_2} & \\
X\ar[rd]_{f} &                                        & Y\ar[ld]^{g} \\
             &    S                                   &
}
$$

Künneth 同构 (\hyperrefext[XVII][XVII.thm5.4.3]{5.4.3})
$$
Rf_{!}(\ZZ/n){\mathop{\otimes}^{\LL}}_{\ZZ/n}Rg_{!}(\ZZ/n)\xrightarrow{\sim}
R(f\times g)_{!}(\ZZ/n)
$$
诱导同构
\begin{equation*}
R^{2d}f_{!}\ZZ/n\otimes
R^{2e}g_{!}\ZZ/n\xrightarrow{\sim} R^{2(d+e)}(f\times
g)_{!}(\ZZ/n)\tag{2.11.1}
\end{equation*}\etiquette[2.11.1]{XVIII.eq2.11.1}

显然 $f\times g$ 满足 $(\ast)_{d+e}$; 有如下相容性:

\begin{enonce*}{命题 2.12}\etiquette[2.12]{XVIII.prop2.12}
下图
$$\vcenter{\etiquette[2.12.1]{XVIII.eq2.12.1}
\xymatrix@R=1.25cm{
R^{2d}f_!\ZZ/n(d) \otimes R^{2e}g_!\ZZ/n (e)
\ar[r]^{\sim} \ar[d]^{\Tr_f\otimes \Tr_g} & R^{2(d+e)}(f\times g)_!
\ZZ/n(d+e) \ar[d]^{\Tr_{f\times g}}\\
\ZZ/n  \ar@{=}[r] & \ZZ/n,
}}\leqno{(2.12.1)}
$$
交换, 其中第一个水平箭头是 \eqhyperref[XVIII.eq2.11.1]{2.11.1} 的扭转形式.\pageoriginale
\end{enonce*}

为验证这一点, 只须利用 (Var 3) 以及在 \hyperrefext[XVII][XVII.thm5.4.3]{5.4.3} 证明中使用的 Künneth 箭头的不对称解释 (\hyperrefext[XVII][XVII.lem5.4.3.5]{5.4.3.5}).

\setcounter{subsection}{12}
\subsection{}\etiquette[2.13]{XVIII.subsec2.13}
设 $f:X\rightarrow S$ 为满足 $(\ast)_{d}$ 的 $S$ 可紧化态射, $\mathcal{A}$ 为 $S$ 上的环层, $n$ 为 $S$ 上可逆且满足 $n\mathcal{A}=0$ 的整数.

对每个 $K\in D(S,\mathcal{A})$, 有 (\hyperrefext[XVII][XVII.thm5.2.6]{5.2.6})
\begin{equation*}
K{\mathop{\otimes}^{\LL}}_{\ZZ/n}Rf_{!}\ZZ/n(d)\xrightarrow{\sim} Rf_{!}(f^{\ast}K(d)){\quoi}.\tag{2.13.1}
\end{equation*}\etiquette[2.13.1]{XVIII.eq2.13.1}

迹态射
$$
\Tr_{f}:R^{2d}f_{!}\ZZ/n(d)\longrightarrow \ZZ/n
$$
也可解释为态射
$$
\Tr_{f}:Rf_{!}\ZZ/n(d)[2d]\longrightarrow \ZZ/n{\quoi}.
$$
将此态射与 $K$ 作张量积, 经由 \eqhyperref[XVIII.eq2.13.1]{2.13.1}, 定义迹态射
\begin{equation*}
\Tr_{f}:Rf_{!}(f^{\ast}K(d)[2d])\longrightarrow K{\quoi}.\tag{2.13.2}
\end{equation*}\etiquette[2.13.2]{XVIII.eq2.13.2}

如将在 \hyperref[XVIII.thm3.2.5]{3.2.5} 中看到的, 下述定理实质上等价于平展上同调中的整体对偶定理.

\begin{enonce*}{定理 2.14}\etiquette[2.14]{XVIII.thm2.14}
设 $f:X\rightarrow S$ 为纯相对维数 $d$ 的光滑可紧化态射, $n$ 为 $S$ 上可逆的整数 $\geq 1$. 对 $X$ 的每个几何点 $x$, 其在 $S$ 中的像为 $s$, 存在 $S$ 中 $s$ 的平展邻域 $V$, 以及 $f^{-1}(V)$ 中 $x$ 的平展邻域 $U$
\begin{equation*}
\vcenter{
\xymatrix@C=1.5cm@R=1.25cm{
U \ar[rd]_{f'_V} \ar[r]^{j} & f^{-1}(V)   \ar[r] \ar[d]^{f_V}& X\ar[d]^{f}\\
                                             &     V           \ar[r]  & S
}}\tag{2.14.1}
\end{equation*}\etiquette[2.14.1]{XVIII.eq2.14.1}%
使得\pageoriginale 态射 \sisi{\emph{(XVII 6.2.7.2)}}{\eqhyperrefext[XVII][XVII.eq6.2.7.4]{6.2.7.4}}
$$
R^{i}f'_{V!}\ZZ/n\longrightarrow R^{i}f_{V!}\ZZ/n
$$
对 $i<2d$ 为零, 且态射
$$
\Tr_{f'_{V}}:R^{2d}f'_{V!}\ZZ/n(d)\longrightarrow \ZZ/n
$$
为同构.
\end{enonce*}

先证明下述引理.

\begin{enonce*}{引理 2.14.2}\etiquette[2.14.2]{XVIII.lem2.14.2}
设 $\mathcal{A}$ 为阿贝尔范畴, $k$ 为整数 $\geq 0$, $(K_{i})_{0\leq i\leq 2k}$ 为 $D^{b}(\mathcal{A})$ 中的对象, 满足对 $p\not\in [0,k]$ 有 $H^{p}(K_{i})=0$, $f_{i}:K_{i}\rightarrow
K_{i+1}$ $(0\leq i\leq 2k-1)$ 为态射, $f$ 为其复合. 若对 $p<k$ 有 $H^{p}(f_{i})=0$, 则在 $D^{b}(\mathcal{A})$ 中存在从复形 $H^{k}(K_{0})[-k]$ (由置于次数 $k$ 的 $H^{k}(K_{0})$ 构成) 到 $K_{2k}$ 的态射 $\varphi$, 使图表
$$
\xymatrix{
K_{0}  \ar[d]\ar[r]^{f} & K_{2k}\\
\sigma_{\geq k}(K_{0})  \ar@{=}[r] & H^{k}(K_{0})[-k] \ar[u]_{\varphi}.
}
$$
交换.
\end{enonce*}

当 $k=0$ 时引理平凡. 对 $k$ 作归纳证明. 对复形 $\sigma_{\geq
1}(K_{i})[1]$ 应用归纳假设, 得到态射 $\varphi':H^{k}(K_{0})[-k]\rightarrow \sigma_{\geq 1}K_{2k-2}$ 使下图交换, 其中 $f'=f_{2k-3}\ldots f_{0}$:
$$
\xymatrix{
K_{0}  \ar[d]\ar[r]^{f'} & \sigma_{\geq 1}K_{2k-2}            \\
\sigma_{\geq k}(K_{0})  \ar@{=}[r] & H^{k}(K_{0})[-k]. \ar[u]_{\varphi'}
}
$$

对每个复形 $L$, 特殊三角
$$
\xymatrix{
                            &\sigma_{\geq 1}(L)\ar[ld]_{+1}  &       \\
  \sigma_{\leq 0}(L) \ar[rr] &                    & L\ar[lu]
}
$$
对每个复形 $M$ 给出\pageoriginale 正合列
$$
\ldots\rightarrow \Hom(M,\sigma_{\leq 0}(L))\rightarrow \Hom
(M,L)\rightarrow \Hom (M,\sigma_{\geq
1}L)\xrightarrow{j}\Ext^{1}(M,\sigma_{\leq 0}L){\quoi}.
$$

特别地, 有图表
$${\def\objectstyle{\scriptscriptstyle}
\xymatrix@C=0.4cm{
\Hom(\sigma_{\geq k}K_{0}, \sigma_{\leq 0} K_{2k-2})\ar[r] \ar[d]&
\overset{\text{\ding{176}}}{\underset{\phantom{\text{\ding{176}}}}
{\Hom}}(\sigma_{\geq k}K_{0},
    K_{2k-2}) \ar[r]\ar[d] &
\overset{\text{\ding{174}}}{\underset{\phantom{\text{\ding{174}}}}
{\Hom}}(\sigma_{\geq
      k}K_{0}, \sigma_{\geq 1} K_{2k-2})\ar[r] \ar[d]&
\overset{\text{\ding{175}}}{\underset{\phantom{\text{\ding{175}}}}
{\Ext}^1}(\sigma_{\geq
      k}K_{0}, \sigma_{\leq 0} K_{2k-2})\ar[d]\\
\overset{\text{\ding{178}}}{\underset{\phantom{\text{\ding{178}}}}
{\Hom}}(K_0,\sigma_{\leq 0} K_{2k-2})\ar[r] &
\overset{\text{\ding{172}}}{\underset{\phantom{\text{\ding{172}}}}
{\Hom}}(K_0,K_{2k-2}) \ar[r] &
\overset{\text{\ding{173}\ding{177}}}{\underset{\phantom{\text{\ding{173}\ding{177}}}}
{\Hom}}(K_0, \sigma_{\geq
  1}K_{2k-2})\ar[r]& \Ext^{1}(K_0, \sigma_{\leq 0} K_{2k-2})
}}\leqno{(2.14.3)}\etiquette[2.14.3]{XVIII.eq2.14.3}
$$
它映到 $K_{2k-1}$ 与 $K_{2k}$ 的类似图表. 根据假设, 位于 \ding{172} 的态射 $f'$ 在 \ding{173} 处的像可提升到 \ding{174}. 将元素 \ding{174} 提升到 \ding{176} 的障碍 \ding{175}, 在 $K_{2k-1}$ 的图表 \eqhyperref[XVIII.eq2.14.3]{2.14.3} 中的像为零. 因而在此新图表中, \ding{174} 可提升到 \ding{176}, 而 \ding{172} 与 \ding{176} 的像之差在 \ding{177} 处的像为零, 因而是某个元素 \ding{178} 的像. 此元素 \ding{178} 在 $K_{2k}$ 的图表 \hyperref[XVIII.eq2.14.3]{2.14.3} 中的像为零; 因而在此新图表中, $f$ 为 \ding{176} 的像, 问题得解.

证明对于给定的 $d$ 值, 断言 \hyperref[XVIII.thm2.14]{2.14} 蕴含下述结果.

\begin{enonce*}{推论 2.14.4}\etiquette[2.14.4]{XVIII.cor2.14.4}
在 \hyperref[XVIII.thm2.14]{2.14} 的假设下, 存在图表 \eqhyperref[XVIII.eq2.14.1]{2.14.1}, 使 \emph{\hyperrefext[XVII][XVII.subsubsec6.2.4]{6.2.4}} 中的态射在 $D^{b}(V,\sheaf{\ZZ}/n)$ 中有分解
$$
\xymatrix@C=.75cm@R=1cm{
Rf'_{V!}\sheaf{\ZZ}/n(d)  \ar[rr] \ar[rd]_{t}&  & Rf_{V!}\sheaf{\ZZ}/n(d)\\
                          &\sheaf{\ZZ}/n[-2d], \ar[ru]_{\varphi}
}
$$
其中 $t$ 为箭头
% SGA-SRC-001151
$$
t:Rf'_{V!}\ZZ/n(d)\rightarrow \sigma_{\geq
2d}Rf'_{V!}\ZZ/n(d)=H^{2d}(Rf'_{V!}\ZZ/n(d))[-2d]\xrightarrow{\Tr}\sheaf{\ZZ}/n[-2d]{\quoi}.
$$
\end{enonce*}

置\pageoriginale $f'_{0}=f$, 归纳构造 $4d$ 个满足 \hyperref[XVIII.thm2.14]{2.14} 的图表 \eqhyperref[XVIII.eq2.14.1]{2.14.1}
$$
\xymatrix@C=2cm@R=1.5cm{
U^{'}_{i+1} \ar[r]\ar[rd]_{f'_{i+1}}& f'^{-1}_{i}(V'_{i+1}) \ar[d]\ar[r]& U'_i \ar[d]^{f'_{i}}\\
            & V'_{i+1}  \ar[r]            & V'_i
}
$$
置 $V=V'_{4d}$ 及 $U_{i}=\fprod{U'_{i}}{V}{V'_{i}}$. 得到图表
$$
\xymatrix@R=1.5cm@C=1cm{
U_{4d}\ar[r] \ar[rrrd]^{f_{4d}}& \cdots \ar[r]& U_1 \ar[r]\ar[rd]^{f_1}& U_0\ar[r]\ar[d]^{f^0} & X \ar[d]^{f}\\
            &              &           & V \ar@{-}[r]  & S ,
}
$$
其中 $U_{0}=f^{-1}(V)$, 使箭头
$$
R^{i}f_{k+1!}\ZZ/n(d)\longrightarrow R^{i}f_{k!}\ZZ/n(d)
$$
为零 (对 $i<2d$), 且箭头
$$
\Tr:R^{2d}f_{4d!}\sheaf{\ZZ}/n(d)\longrightarrow \sheaf{\ZZ}/n
$$
为同构. 由 \hyperref[XVIII.lem2.14.2]{2.14.2}, 推论对 $U=U_{4d}$ 成立.

现在已\sisi{准备就绪}{准备就绪}, 对 $d$ 作归纳证明 \hyperref[XVIII.thm2.14]{2.14}. 情形 $d=0$ 平凡, 情形 $d=1$ 恰为 \hyperref[XVIII.lemfond1.6.9]{1.6.9}. 因而可以假设 $d\geq 2$.

由于问题在 $X$ 与 $S$ 上是局部的, 可以假设 $f$ 有分解 $f=gh$, 其中 $g$ 与 $h$ 光滑且可紧化, 纯相对维数为 $d'$ 及 $d''$, 并且 $d'$、$d''<d=d'+d''$. 对 $g$ 以及由 $h$ 经基变换所得的态射应用 \hyperref[XVIII.cor2.14.4]{2.14.4}, 得到图表
$$
\xymatrix@C=2.2cm@R=1.5cm{
  U  \ar[r]^{u} \ar[rd]^{h'_0} & f^{-1}W \ar[d]^{h_0}       \ar[r] & X \ar[d]^{h} \\
V  \ar[r]^{v} \ar[rd]^{g'_0} & g^{-1}W \ar[r] \ar[d]^{g_0}       & Y \ar[d]^{g}\\
               &    W   \ar[r]               & S,
}
$$
使得\pageoriginale 态射 (XVII \hyperref[XVII.eq6.2.7.3]{6.2.7.3}) 有分解
\begin{align*}
 Rh'_{0!}\ZZ/n(d'') & \xrightarrow{\Tr}
\ZZ/n[-2d'']\xrightarrow{\varphi} Rh_{0!}\ZZ/n(d'')\\
Rg'_{0!}\ZZ/n(d') & \xrightarrow{\Tr}\ZZ/n[-2d']\xrightarrow{\Psi} Rg_{0!}\ZZ/n(d')
\end{align*}
且 $\Tr_{h'_0}$ 与 $\Tr_{g'_0}$ 为同构.

设 $h_{1}$ 与 $h'_{1}$ 为由 $h_{0}$ 与 $h'_0$ 经基变换到 $V$ 所得的态射, 并设 $f'=g'_{0}h'_{1}$、$f_{W}=g_{0}h_{0}$.

下图交换:
$$
\xymatrix{
  Rf'_{!}\ZZ/n(d) \ar[rr]\ar@{=}[d] & & Rf_{W!}\ZZ/n(d) \ar@{=}[d]\\
Rg'_{0!} Rh'_1\ZZ/n(d) \ar[r]& Rg_{0!} Rh'_{0!}\ZZ/n(d) \ar[r]& Rg_{0!}
Rh_{0!}\ZZ/n(d),
}
$$
图表
$$
\xymatrix@C=2.1cm@R=1cm{
  Rg'_{0!} Rh'_{1!} \ZZ/n(d)   \ar@{=}[d] \ar@<-5ex>@/_2pc/[ddd]& &\\
Rg'_{0!} v^{\ast} Rh'_{0!} \ZZ/n(d)
 \ar[dd] \ar[r]_{Rg'_{0!}v^{\ast}(t)(d')}&
Rg'_{0!}\ZZ/n(d')[-2d''] \ar[d]^{t[-2d'']}& \\
   & \ZZ/n[-2d] \ar[d]^{\Psi[-2d'']}\\
Rg_{0!} Rh'_{0!} \ZZ/n(d) \ar[r]_{Rg_{0!}(t)(d')} \ar@/_2.2pc/[rr]& Rg_{0!}\ZZ/n(d')[-2d''] \ar[r]_{Rg_{0!}(\varphi)(d')}& Rg_{0!}Rh_{0!}\ZZ/n(d)
}
$$
也交换. 因而对于\pageoriginale 图表
$$
\xymatrix@C=1.5cm@R=1.25cm{
U_1 \ar[r]\ar[rd]^{f'} & f^{-1}W \ar[r]\ar[d]^{f_W}& X\ar[d]^{f}\\
                       &  W \ar@{-}[r]     & S,
}
$$
从 $Rf'_{!}\ZZ/n(d)$ 到 $Rf_{W!}\ZZ/n(d)$ 的典范箭头经由 $\ZZ/n[-2d]$ 分解; 因而从 $R^{i}f'_{!}\ZZ/n(d)$ 到 $R^{i}f_{W!}\ZZ/n(d)$ 的箭头对 $i<2d$ 为零. 最后, $\Tr_{f'}$ 为同构, 因为它是同构 $\Tr_{g'_{0}}$ 与 $R^{2d'}g'_{0!}\Tr_{h'_{1}}$ 的复合.

至此完成了 \hyperref[XVIII.thm2.14]{2.14} 的证明.


